\documentclass[a4paper,12pt]{article}

% ---------- Кодировка и язык ----------
\usepackage[T2A]{fontenc}
\usepackage[utf8]{inputenc}
\usepackage[russian]{babel}

% ---------- Математика ----------
\usepackage{amsmath,amssymb,mathtools}
\usepackage{icomma}

% ---------- Шрифт и типографика ----------
\usepackage{helvet}
\renewcommand{\familydefault}{\sfdefault}
\usepackage{microtype}
\usepackage{setspace}
\usepackage{parskip}

% ---------- Страница ----------
\usepackage[
  a4paper,
  left=18mm,
  right=18mm,
  top=18mm,
  bottom=20mm
]{geometry}

% ---------- Таблицы и списки ----------
\usepackage{array}
\usepackage{tabularx}
\usepackage{booktabs}
\usepackage{enumitem}

% ---------- Графика ----------
\usepackage{tikz}
\usetikzlibrary{calc}

% ---------- Цвет ----------
\usepackage{xcolor}
\definecolor{accent}{HTML}{008080}
\definecolor{darkaccent}{HTML}{004D4D}
\definecolor{textgray}{HTML}{333333}
\definecolor{softgray}{HTML}{F2F4F4}

% ---------- Ссылки ----------
\usepackage{hyperref}
\hypersetup{
  colorlinks=true,
  linkcolor=darkaccent,
  urlcolor=darkaccent
}

% ---------- Колонтитулы ----------
\usepackage{fancyhdr}
\setlength{\headheight}{15pt}
\pagestyle{fancy}
\fancyhf{}
\fancyhead[L]{\small\color{textgray}Чек-лист}
\fancyhead[R]{\small\color{textgray}Тимофей Васильченко}
\fancyfoot[C]{\small\thepage}
\renewcommand{\headrulewidth}{0.4pt}
\renewcommand{\headrule}{%
  \hbox to\headwidth{\color{accent!55}\leaders\hrule height \headrulewidth\hfill}
}

% ---------- Общие настройки ----------
\onehalfspacing
\setlength{\parindent}{0pt}
\setlength{\parskip}{0.45em}
\setlength{\tabcolsep}{5pt}
\renewcommand{\arraystretch}{1.3}
\emergencystretch=2em
\raggedbottom

\setlist[itemize]{
  leftmargin=1.6em,
  itemsep=0.25em,
  topsep=0.35em
}

\setlist[enumerate]{
  leftmargin=1.8em,
  itemsep=0.55em,
  topsep=0.4em
}

\newcommand{\checksection}[1]{%
  \section*{\color{darkaccent}#1}
  \vspace{-0.6em}
  {\color{accent!70}\rule{\linewidth}{0.6pt}}
  \vspace{0.35em}
}

\newcommand{\keyidea}[1]{%
  \textcolor{darkaccent}{\textbf{#1}}
}

\begin{document}

% ============================================================
% ТИТУЛЬНЫЙ ЛИСТ
% ============================================================

\begin{titlepage}
\thispagestyle{empty}

\begin{center}

\vspace*{22mm}

{\Large\bfseries\color{darkaccent} ЧЕК-ЛИСТ}

\vspace{12mm}

{\Huge\bfseries
Текстовые задачи}

\vspace{5mm}

{\LARGE\bfseries\color{accent}
табличный метод, движение,\\
производительность и растворы}

\vspace{10mm}

{\Large
Работа с квадратным уравнением\\
и «проблемным» дискриминантом}

\vspace{22mm}

{\large
Лёвин Артём Александрович\\[2mm]
\textbf{эксклюзивно для Тимофея Васильченко}
}

\vfill

{\large \today}

\vspace{18mm}

\end{center}

% Сдержанная кривая Безье внизу титульного листа
\begin{tikzpicture}[remember picture,overlay]
  \draw[
    accent!55,
    line width=1.1pt
  ]
  ([xshift=18mm,yshift=13mm]current page.south west)
  .. controls
  ([xshift=55mm,yshift=27mm]current page.south west)
  and
  ([xshift=-58mm,yshift=8mm]current page.south east)
  ..
  ([xshift=-18mm,yshift=17mm]current page.south east);
\end{tikzpicture}

\end{titlepage}

% ============================================================
% ОСНОВНОЙ ЧЕК-ЛИСТ
% ============================================================

\checksection{1. Универсальная схема решения текстовой задачи}

Для задач на движение, производительность и растворы полезно придерживаться одного и того же каркаса.

\begin{enumerate}
  \item \keyidea{Определите величины}, которые описывают ситуацию.
  \item \keyidea{Выберите неизвестную.} Обычно удобно обозначать через \(x\) именно ту величину, которую требуется найти.
  \item \keyidea{Запишите единицы измерения} и приведите их к единой системе.
  \item \keyidea{Составьте таблицу} и сначала внесите в неё данные из условия.
  \item \keyidea{Заполните вычисляемые ячейки} с помощью основной формулы задачи.
  \item \keyidea{Переведите словесную связь в уравнение} или систему уравнений.
  \item \keyidea{Запишите ограничения} на неизвестные и знаменатели.
  \item \keyidea{Решите уравнение} удобным способом.
  \item \keyidea{Отберите допустимый ответ} по смыслу задачи.
\end{enumerate}

\begin{table}[ht]
\centering
\caption{Три основные модели}
\begin{tabularx}{\linewidth}{@{}p{3.1cm}p{5.1cm}X@{}}
\toprule
\textbf{Тип задачи} & \textbf{Величины} & \textbf{Основная связь} \\
\midrule
Движение
&
скорость \(v\), время \(t\), путь \(s\)
&
\[
s=vt,\qquad
t=\frac{s}{v},\qquad
v=\frac{s}{t}
\]
\\
Производительность
&
производительность \(p\), время \(t\), работа \(W\)
&
\[
W=pt,\qquad
t=\frac{W}{p}
\]
\\
Растворы
&
масса раствора \(m\), концентрация \(c\), масса вещества \(q\)
&
\[
q=cm
\]
\\
\bottomrule
\end{tabularx}
\end{table}

\keyidea{Главная идея:} таблица нужна не сама по себе. Она организует данные так, чтобы нужное уравнение становилось почти очевидным.

% ------------------------------------------------------------

\checksection{2. Задачи на движение}

В таблице движения удобно использовать три столбца:

\[
\boxed{\text{скорость}}
\qquad
\boxed{\text{время}}
\qquad
\boxed{\text{расстояние}}.
\]

Если два объекта проходят один и тот же путь, расстояния в соответствующих строках таблицы совпадают.

\subsection*{\color{darkaccent}Пример}

Первый теплоход проходит \(323\) км со скоростью \(x\) км/ч. Второй движется на \(2\) км/ч быстрее и отправляется на \(2\) часа позже. Оба прибывают одновременно.

\begin{table}[ht]
\centering
\begin{tabularx}{0.88\linewidth}{@{}lccc@{}}
\toprule
& \textbf{Скорость} & \textbf{Время} & \textbf{Путь} \\
\midrule
Первый
&
\(x\)
&
\(\dfrac{323}{x}\)
&
\(323\)
\\[2mm]
Второй
&
\(x+2\)
&
\(\dfrac{323}{x+2}\)
&
\(323\)
\\
\bottomrule
\end{tabularx}
\end{table}

Поскольку первый теплоход отправился раньше, время его движения на \(2\) часа больше:

\[
\frac{323}{x}-\frac{323}{x+2}=2.
\]

По смыслу задачи

\[
x>0.
\]

Следовательно, знаменатели также отличны от нуля.

Приводим левую часть к одной дроби:

\[
\frac{323(x+2)-323x}{x(x+2)}=2.
\]

В числителе происходит существенное сокращение:

\[
\frac{646}{x(x+2)}=2.
\]

Отсюда

\[
646=2x(x+2),
\]

\[
323=x^2+2x,
\]

\[
x^2+2x-323=0.
\]

Получаем квадратное уравнение, в котором требуется аккуратно вычислить дискриминант.

\[
D=2^2-4\cdot1\cdot(-323)
 =4+1292
 =1296.
\]

\[
\sqrt D=36.
\]

Тогда

\[
x_{1,2}=\frac{-2\pm36}{2},
\]

\[
x_1=17,\qquad x_2=-19.
\]

Отрицательная скорость не соответствует смыслу задачи, поэтому

\[
\boxed{x=17\text{ км/ч}}.
\]

\subsection*{\color{darkaccent}Как правильно составлять разность времён}

Фраза «выехал раньше» сама по себе ещё не задаёт уравнение. Нужно восстановить хронологию.

Например:

\begin{itemize}
  \item второй выехал на \(2\) часа позже и прибыли одновременно:
  \[
  t_1-t_2=2;
  \]

  \item второй выехал на \(1\) час позже, но приехал на \(1\) час раньше:
  \[
  t_1-t_2=2.
  \]
\end{itemize}

Поэтому сначала определяйте, \keyidea{на сколько различается фактическое время движения}, и только затем составляйте уравнение.

% ------------------------------------------------------------

\checksection{3. Алгебраические дроби: работаем без лишних вычислений}

В текстовых задачах часто появляются выражения вида

\[
\frac{A}{p(x)}-\frac{B}{q(x)}
\qquad\text{или}\qquad
\frac{A}{p(x)}+\frac{B}{q(x)}.
\]

Перед преобразованиями необходимо записать

\[
p(x)\ne0,\qquad q(x)\ne0.
\]

Общий знаменатель можно записать сразу:

\[
\frac{A}{p}-\frac{B}{q}
=
\frac{Aq-Bp}{pq}.
\]

При этом не следует торопиться перемножать большие числа. Сначала ищите общие множители и сокращения.

\subsection*{\color{darkaccent}Пример}

\[
\frac{x}{28}+\frac{x}{22}=25.
\]

Общий знаменатель можно не вычислять отдельно:

\[
\frac{22x+28x}{28\cdot22}=25,
\]

\[
\frac{50x}{28\cdot22}=25.
\]

Отсюда

\[
50x=25\cdot28\cdot22.
\]

Сокращаем до умножения:

\[
x=\frac{28\cdot22}{2}=14\cdot22=308.
\]

\[
\boxed{x=308}.
\]

\keyidea{Важно:} сокращать можно только множители. Нельзя «сокращать» отдельные слагаемые внутри суммы или разности.

% ------------------------------------------------------------

\checksection{4. Как извлекать корень из «неудобного» дискриминанта}

После правильного составления текстовой задачи основная вычислительная трудность нередко возникает на этапе

\[
x_{1,2}=\frac{-b\pm\sqrt D}{2a}.
\]

Полезно владеть несколькими способами точного нахождения \(\sqrt D\).

\subsection*{\color{darkaccent}Метод 1. Разложение на множители}

Если число удобно раскладывается, используем разложение степеней.

Например,

\[
1296=2^4\cdot3^4.
\]

Тогда

\[
\sqrt{1296}
=
\sqrt{2^4\cdot3^4}
=
2^2\cdot3^2
=
4\cdot9
=
36.
\]

То есть

\[
\boxed{\sqrt{1296}=36}.
\]

Общее правило:

\[
\sqrt{a^{2m}}=|a^m|.
\]

Для положительных чисел, возникающих в разложении дискриминанта, модуль обычно не создаёт дополнительной сложности.

\subsection*{\color{darkaccent}Метод 2. Ближайшие квадраты и последняя цифра}

Сначала определяем диапазон для корня.

Например,

\[
70^2=4900,\qquad
80^2=6400.
\]

Если

\[
4900<D<6400,
\]

то

\[
70<\sqrt D<80.
\]

Следовательно, если \(D\) является полным квадратом, корень имеет вид \(7\Box\).

Далее используем последнюю цифру числа.

\begin{table}[ht]
\centering
\caption{Последняя цифра полного квадрата}
\begin{tabular}{@{}cc@{}}
\toprule
\textbf{Последняя цифра квадрата} &
\textbf{Возможная последняя цифра корня}
\\
\midrule
\(0\) & \(0\) \\
\(1\) & \(1\) или \(9\) \\
\(4\) & \(2\) или \(8\) \\
\(5\) & \(5\) \\
\(6\) & \(4\) или \(6\) \\
\(9\) & \(3\) или \(7\) \\
\bottomrule
\end{tabular}
\end{table}

Число, оканчивающееся на \(2,3,7\) или \(8\), не может быть квадратом целого числа.

\subsection*{\color{darkaccent}Пример: \(\sqrt{5329}\)}

Имеем

\[
70^2<5329<80^2.
\]

Следовательно,

\[
\sqrt{5329}=7\Box.
\]

Число \(5329\) оканчивается на \(9\), поэтому возможны только окончания \(3\) и \(7\):

\[
73\qquad\text{или}\qquad77.
\]

Проверяем более подходящего кандидата:

\[
73^2=(70+3)^2
=4900+420+9
=5329.
\]

Значит,

\[
\boxed{\sqrt{5329}=73}.
\]

Такой подход особенно удобен, когда разложение числа перебором делителей оказывается слишком долгим.

\subsection*{\color{darkaccent}Ещё пример: \(\sqrt{5184}\)}

\[
70^2<5184<80^2.
\]

Последняя цифра равна \(4\), поэтому рассматриваем

\[
72\qquad\text{и}\qquad78.
\]

\[
72^2=(70+2)^2
=4900+280+4
=5184.
\]

Следовательно,

\[
\boxed{\sqrt{5184}=72}.
\]

\keyidea{Важно:} последняя цифра только сокращает количество кандидатов. Она не доказывает, что найденное число действительно является корнем. Кандидата нужно проверить возведением в квадрат.

% ------------------------------------------------------------

\checksection{5. Огромный дискриминант: корень рядом с \(\lvert b\rvert\)}

Пусть

\[
D=b^2-4ac.
\]

Если \(4ac>0\) и величина \(4ac\) сравнительно мала относительно \(b^2\), то

\[
D<b^2,
\]

поэтому

\[
\sqrt D<|b|,
\]

причём корень может находиться совсем близко к \(|b|\).

Это позволяет искать кандидата около \(|b|\), вместо разложения огромного числа на простые множители.

\subsection*{\color{darkaccent}Пример}

Решим

\[
100x^2-5002x+100=0.
\]

Здесь

\[
a=100,\qquad b=-5002,\qquad c=100.
\]

Дискриминант:

\[
D=(-5002)^2-4\cdot100\cdot100.
\]

\[
5002^2
=
(5000+2)^2
=
25\,020\,004.
\]

Поэтому

\[
D
=
25\,020\,004-40\,000
=
24\,980\,004.
\]

Разложение такого числа не выглядит привлекательным.

Но

\[
D=5002^2-40\,000,
\]

поэтому \(\sqrt D\) должен быть немного меньше \(5002\).

Число \(D\) оканчивается на \(4\), следовательно, целый корень должен оканчиваться на \(2\) или \(8\).

Ближайший подходящий кандидат ниже \(5002\):

\[
4998.
\]

Проверяем:

\[
4998^2
=
(5000-2)^2
=
25\,000\,000-20\,000+4
=
24\,980\,004.
\]

Значит,

\[
\boxed{\sqrt D=4998}.
\]

Теперь

\[
x_{1,2}
=
\frac{5002\pm4998}{200}.
\]

Получаем

\[
x_1=\frac{10000}{200}=50,
\]

\[
x_2=\frac{4}{200}=\frac1{50}.
\]

Итак,

\[
\boxed{x=50\quad\text{или}\quad x=\frac1{50}}.
\]

\keyidea{Принцип метода:}

\[
\boxed{
D=b^2-\text{сравнительно небольшое число}
\quad\Longrightarrow\quad
\sqrt D\text{ ищем около }|b|
}
\]

Оценка даёт кандидата, а точность обязательно подтверждается возведением в квадрат.

% ------------------------------------------------------------

\checksection{6. Как выбрать способ вычисления \(\sqrt D\)}

\begin{table}[ht]
\centering
\caption{Рациональный выбор метода}
\begin{tabularx}{\linewidth}{@{}p{4.2cm}X@{}}
\toprule
\textbf{Ситуация} & \textbf{Что делать} \\
\midrule
Число удобно делится на небольшие простые множители
&
Разложить на множители и вынести пары из-под корня.
\\
Число похоже на полный квадрат и находится в понятном диапазоне
&
Зажать его между ближайшими квадратами и использовать последнюю цифру.
\\
Разложение выглядит долгим
&
Сначала оценить диапазон корня; часто это быстрее перебора делителей.
\\
\(D=b^2-4ac\), причём \(4ac\) мало по сравнению с \(b^2\)
&
Искать \(\sqrt D\) около \(|b|\), затем проверить кандидата квадратом.
\\
\bottomrule
\end{tabularx}
\end{table}

Метод ближайших квадратов \keyidea{не ограничен числами меньше \(10\,000\)}. Для больших чисел он также корректен, если удаётся быстро подобрать удобные ориентиры.

% ------------------------------------------------------------

\checksection{7. Задачи на производительность}

Производительность --- это объём выполненной работы за единицу времени.

Например:

\begin{itemize}
  \item литры в минуту;
  \item детали в час;
  \item вопросы в час;
  \item страницы в минуту.
\end{itemize}

Основная формула:

\[
W=pt.
\]

Отсюда

\[
t=\frac{W}{p}.
\]

Если первая труба имеет производительность \(x\) л/мин, а вторая пропускает на \(5\) л/мин больше, то

\[
p_1=x,\qquad p_2=x+5.
\]

Если обе трубы должны пропустить один и тот же объём \(W\), то

\[
t_1=\frac{W}{x},
\qquad
t_2=\frac{W}{x+5}.
\]

Если первая выполняет работу на \(5\) минут дольше:

\[
\frac{W}{x}-\frac{W}{x+5}=5.
\]

\keyidea{Единицы измерения определяются производительностью.}

Если

\[
p=6{,}7\ \text{вопроса/ч},
\]

то время в формуле

\[
t=\frac{W}{p}
\]

должно быть выражено в часах.

Например,

\[
30\text{ мин}=\frac12\text{ ч}.
\]

% ------------------------------------------------------------

\checksection{8. Задачи на растворы и смеси}

Для каждого раствора используются три величины:

\[
\boxed{\text{масса раствора}}
\qquad
\boxed{\text{концентрация}}
\qquad
\boxed{\text{масса вещества}}.
\]

Основная связь:

\[
q=cm.
\]

Здесь

\[
m=\text{масса раствора},
\qquad
c=\text{концентрация},
\qquad
q=\text{масса растворённого вещества}.
\]

Процентную концентрацию переводим в долю:

\[
45\%=0{,}45,
\qquad
97\%=0{,}97,
\qquad
50\%=0{,}5.
\]

Для чистой воды

\[
c=0,
\]

поскольку растворённого вещества в ней нет.

Для чистого вещества

\[
c=1.
\]

\keyidea{Концентрации складывать нельзя.}

Складываются:

\[
\text{массы растворов}
\]

и

\[
\text{массы растворённого вещества}.
\]

% ------------------------------------------------------------

\subsection*{\color{darkaccent}Пример с двумя экспериментами}

Пусть имеется:

\begin{itemize}
  \item \(x\) кг \(45\%\)-го раствора;
  \item \(y\) кг \(97\%\)-го раствора.
\end{itemize}

В первом случае добавляют \(10\) кг чистой воды и получают \(62\%\)-й раствор.

\begin{table}[ht]
\centering
\begin{tabularx}{0.94\linewidth}{@{}Xccc@{}}
\toprule
\textbf{Компонент}
&
\textbf{Масса}
&
\textbf{Концентрация}
&
\textbf{Масса вещества}
\\
\midrule
Первый раствор
&
\(x\)
&
\(0{,}45\)
&
\(0{,}45x\)
\\
Второй раствор
&
\(y\)
&
\(0{,}97\)
&
\(0{,}97y\)
\\
Вода
&
\(10\)
&
\(0\)
&
\(0\)
\\
После смешивания
&
\(x+y+10\)
&
\(0{,}62\)
&
\(0{,}62(x+y+10)\)
\\
\bottomrule
\end{tabularx}
\end{table}

Получаем

\[
0{,}45x+0{,}97y
=
0{,}62(x+y+10).
\tag{1}
\]

Во втором случае вместо чистой воды добавляют \(10\) кг \(50\%\)-го раствора и получают концентрацию \(72\%\).

Тогда

\[
0{,}45x+0{,}97y+0{,}5\cdot10
=
0{,}72(x+y+10).
\tag{2}
\]

Получаем систему

\[
\begin{cases}
0{,}45x+0{,}97y
=
0{,}62(x+y+10),\\[1mm]
0{,}45x+0{,}97y+5
=
0{,}72(x+y+10).
\end{cases}
\]

Вычитая первое уравнение из второго:

\[
5=0{,}1(x+y+10),
\]

\[
x+y+10=50,
\]

\[
x+y=40.
\]

Из первого уравнения

\[
0{,}45x+0{,}97y=31.
\]

Решение системы:

\[
\boxed{x=15,\qquad y=25}.
\]

Здесь необходимы два независимых уравнения, поскольку имеются две неизвестные величины.

% ------------------------------------------------------------

\checksection{9. Что особенно важно не перепутать}

\begin{itemize}
  \item Скорость, производительность и масса раствора по смыслу задачи должны быть положительными.
  
  \item Перед умножением рационального уравнения на знаменатели нужно проверить, при каких значениях они обращаются в ноль.
  
  \item Если получен отрицательный корень квадратного уравнения, его нельзя автоматически отбрасывать. Сначала нужно проверить смысл исходной задачи.
  
  \item В задачах на движение разность времён определяется всей хронологией отправления и прибытия.
  
  \item Минуты и часы нельзя смешивать в одном уравнении без предварительного перевода.
  
  \item Концентрация \(45\%\) в вычислениях записывается как \(0{,}45\).
  
  \item Концентрации компонентов не складываются. Складываются массы и массы растворённого вещества.
  
  \item Количество уравнений должно быть достаточным для определения неизвестных; уравнения при этом должны быть независимыми.
  
  \item Последняя цифра дискриминанта лишь помогает отобрать кандидатов для \(\sqrt D\).
  
  \item Найденный подбором корень из большого числа обязательно проверяется возведением в квадрат.
  
  \item При больших коэффициентах сначала ищите сокращения. Перемножать крупные числа следует только тогда, когда это действительно необходимо.
\end{itemize}

% ------------------------------------------------------------

\checksection{10. Короткий алгоритм перед экзаменом}

Перед тем как решать текстовую задачу, проверьте:

\begin{enumerate}
  \item Что требуется найти?
  \item Какие три величины образуют основную модель?
  \item Какие данные можно сразу внести в таблицу?
  \item В одинаковых ли единицах записаны величины?
  \item Какая формула заполняет оставшийся столбец?
  \item Какая фраза из условия превращается в уравнение?
  \item Какие ограничения имеет \(x\)?
  \item Можно ли упростить дроби до перемножения больших чисел?
  \item Как рациональнее найти \(\sqrt D\)?
  \item Подходит ли найденный корень по смыслу задачи?
\end{enumerate}

% ============================================================
% ТРЕНИРОВОЧНЫЙ БЛОК
% ============================================================

\clearpage

\thispagestyle{fancy}

\begin{center}
{\LARGE\bfseries\color{darkaccent}Тренировочный блок}

\vspace{2mm}

{\large Путь А}
\end{center}

\vspace{3mm}

\textbf{Указание.}
Для текстовых задач составляйте таблицу, записывайте уравнение и только затем переходите к вычислениям. Если возникает квадратное уравнение, выбирайте рациональный способ вычисления \(\sqrt D\).

\subsection*{\color{darkaccent}Средний уровень}

\begin{enumerate}[label=\textbf{\arabic*.}]

\item
Первый теплоход проходит расстояние \(240\) км с постоянной скоростью. Второй теплоход движется на \(4\) км/ч быстрее и отправляется на \(2\) часа позже. В пункт назначения они прибывают одновременно.

Найдите скорость первого теплохода.

\item
Первая труба пропускает \(180\) л воды за некоторое время. Вторая труба пропускает такой же объём воды, имея производительность на \(6\) л/мин больше. Первая труба работает на \(5\) минут дольше.

Найдите производительность каждой трубы.

\item
Смешали \(20\%\)-й и \(50\%\)-й растворы и получили \(12\) кг \(35\%\)-го раствора.

Сколько килограммов каждого исходного раствора было взято?

\end{enumerate}

\subsection*{\color{darkaccent}Выше среднего}

\begin{enumerate}[label=\textbf{\arabic*.},resume]

\item
Первый автомобиль проходит \(280\) км с постоянной скоростью. Второй автомобиль едет на \(7\) км/ч быстрее, отправляется на \(1\) час позже первого, но прибывает на \(1\) час раньше.

Найдите скорости обоих автомобилей.

\item
Два насоса должны перекачать по \(336\) л воды. Производительность второго насоса на \(8\) л/мин больше производительности первого, поэтому второй выполняет работу на \(7\) минут быстрее.

Найдите производительность каждого насоса.

\item
Два автомата поочерёдно обработали по одной одинаковой партии деталей. Первый обрабатывает \(28\) деталей в час, второй --- \(22\) детали в час. Суммарное время обработки двух партий составило \(25\) часов.

Сколько деталей было в каждой партии?

\item
Имеются неизвестные массы \(30\%\)-го и \(70\%\)-го растворов.

Если к ним добавить \(10\) кг чистой воды, получится раствор концентрации \(45\%\).

Если вместо воды добавить \(10\) кг \(60\%\)-го раствора, получится раствор концентрации \(55\%\).

Найдите массы двух исходных растворов.

\item
Решите квадратное уравнение

\[
x^2-80x+304=0.
\]

Корень из дискриминанта найдите методом ближайших квадратов и последней цифры.

\item
Решите уравнение

\[
100x^2-5002x+100=0,
\]

не раскладывая дискриминант на простые множители. Используйте близость \(\sqrt D\) к \(|b|\).

\end{enumerate}

\subsection*{\color{darkaccent}Сложная задача}

\begin{enumerate}[label=\textbf{\arabic*.},resume]

\item
Первый теплоход проходит \(1326\) км с постоянной скоростью. Второй движется на \(5\) км/ч быстрее, выходит из пункта отправления на \(2\) часа позже первого и прибывает в пункт назначения на \(3\) часа раньше.

Найдите скорости обоих теплоходов.

При вычислении дискриминанта не используйте перебор делителей. Определите \(\sqrt D\) через ближайшие квадраты и анализ последней цифры.

\end{enumerate}

% ============================================================
% ОТВЕТЫ
% ============================================================

\clearpage

\begin{center}
{\LARGE\bfseries\color{darkaccent}Ответы к тренировочному блоку}
\end{center}

\vspace{4mm}

\begin{table}[ht]
\centering
\caption{Ключевые уравнения и ответы}
\small
\begin{tabularx}{\linewidth}{@{}c X >{\raggedright\arraybackslash}p{4.2cm}@{}}
\toprule
\textbf{№}
&
\textbf{Ключевая модель}
&
\textbf{Ответ}
\\
\midrule

1
&
\(\displaystyle
\frac{240}{x}-\frac{240}{x+4}=2
\)
&
\(20\) км/ч
\\[3mm]

2
&
\(\displaystyle
\frac{180}{x}-\frac{180}{x+6}=5
\)
&
\(12\) л/мин и \(18\) л/мин
\\[3mm]

3
&
\(\displaystyle
0{,}2x+0{,}5(12-x)=0{,}35\cdot12
\)
&
\(6\) кг и \(6\) кг
\\[3mm]

4
&
\(\displaystyle
\frac{280}{x}-\frac{280}{x+7}=2
\)
&
\(28\) км/ч и \(35\) км/ч
\\[3mm]

5
&
\(\displaystyle
\frac{336}{x}-\frac{336}{x+8}=7
\)
&
\(16\) л/мин и \(24\) л/мин
\\[3mm]

6
&
\(\displaystyle
\frac{x}{28}+\frac{x}{22}=25
\)
&
\(308\) деталей
\\[3mm]

7
&
\(\displaystyle
\begin{cases}
0{,}3x+0{,}7y=0{,}45(x+y+10),\\
0{,}3x+0{,}7y+6=0{,}55(x+y+10)
\end{cases}
\)
&
\(20\) кг \(30\%\)-го и \(30\) кг \(70\%\)-го раствора
\\[5mm]

8
&
\(\displaystyle
D=5184,\qquad \sqrt D=72
\)
&
\(x=4,\ 76\)
\\[3mm]

9
&
\(\displaystyle
D=24\,980\,004=4998^2
\)
&
\(\displaystyle x=50,\ \frac1{50}\)
\\[3mm]

10
&
\(\displaystyle
\frac{1326}{x}-\frac{1326}{x+5}=5,
\quad
D=5329=73^2
\)
&
\(34\) км/ч и \(39\) км/ч
\\

\bottomrule
\end{tabularx}
\end{table}

\end{document}